% Luc Destombe --- licence CC BY-NC-SA 4.0
% https://creativecommons.org/licenses/by-nc-sa/4.0/deed.fr
% Fichier autonome : compiler avec pdflatex, deux fois (signets, nombre de pages).
\def\lyceeversion{1.1}
\documentclass[11pt,a4paper]{article}
\makeatletter
\newif\ifeleve
\elevefalse

\newif\iflycee@palatino
\lycee@palatinotrue

\RequirePackage[utf8]{inputenc}
\RequirePackage[T1]{fontenc}
\RequirePackage[french]{babel}
\RequirePackage{etoolbox}

\newcommand{\ShorthandsOff}{\@ifundefined{shorthandoff}{}{\shorthandoff{;:!?}}}
\newcommand{\ShorthandsOn}{\@ifundefined{shorthandon}{}{\shorthandon{;:!?}}}
\ShorthandsOff
\AtBeginDocument{\ShorthandsOn}
\AtBeginEnvironment{tikzpicture}{\ShorthandsOff}

\RequirePackage{geometry}
\RequirePackage{fancyhdr}
\RequirePackage{lastpage}
\RequirePackage{multicol}
\RequirePackage{array}
\RequirePackage{enumitem}
\RequirePackage{xcolor}
\RequirePackage{needspace}

\RequirePackage{amsmath,amssymb}
\RequirePackage{mathpazo}
\renewcommand{\ttdefault}{lmtt}
\DeclareMathSizes{10.95}{10.95}{8}{5.5}
\begingroup\catcode`\_=\active \catcode`\^=\active
\gdef_#1{\sb{\scriptscriptstyle #1}}
\gdef^#1{\sp{\scriptscriptstyle\lycee@moinsepais #1}}
\endgroup
\newcommand\lycee@moins{\mathbin{%
  \setbox\z@\hbox{$\m@th\scriptscriptstyle\lycee@moinsorig$}%
  \dimen@=.55\wd\z@ \advance\dimen@-.5pt
  \hbox to.75\wd\z@{\hss$\m@th\scriptscriptstyle\vcenter{\hbox{%
    \tikz\draw[line width=.5pt,line cap=round](0,0)--(\the\dimen@,0);}}$\hss}}}
\begingroup\lccode`\~=`\- \lowercase{\endgroup\let~}\lycee@moins
\newcommand\lycee@moinsepais{\mathcode`\-="8000 }
\AtBeginDocument{\mathchardef\lycee@moinsorig=\mathcode`\-\relax}
\AtBeginDocument{\catcode`\_=12 \mathcode`\_="8000
  \catcode`\^=12 \mathcode`\^="8000 }
\newcommand\lycee@exposants{%
  \fontdimen13\textfont2=.48\fontdimen6\textfont2
  \fontdimen14\textfont2=.48\fontdimen6\textfont2
  \fontdimen15\textfont2=.48\fontdimen6\textfont2
  \fontdimen13\scriptfont2=.48\fontdimen6\scriptfont2
  \fontdimen14\scriptfont2=.48\fontdimen6\scriptfont2
  \fontdimen15\scriptfont2=.48\fontdimen6\scriptfont2}
\every@math@size\expandafter{\the\every@math@size\lycee@exposants}
\newlength{\retraitcalculs}
\setlength{\retraitcalculs}{2em}

\RequirePackage{tikz}
\usetikzlibrary{arrows.meta,calc,babel,shapes.misc}
\RequirePackage[most]{tcolorbox}
\tcbuselibrary{skins,breakable}

\definecolor{coulDef}{HTML}{1F4E79}
\definecolor{coulProp}{HTML}{7030A0}
\definecolor{coulRem}{HTML}{C55A11}
\definecolor{coulEx}{HTML}{2E7D32}
\definecolor{coulExo}{HTML}{B03A2E}
\definecolor{coulPart}{HTML}{1F4E79}
\definecolor{coulMeth}{HTML}{1B6E4F}
\definecolor{coulCourbe}{HTML}{0B5ED7}
\definecolor{coulCourbeB}{HTML}{C00000}
\definecolor{coulCourbeC}{HTML}{2E7D32}
\definecolor{coulGrille}{HTML}{8C8C8C}
\definecolor{coulRep}{HTML}{1B6E4F}
\definecolor{coulBar}{HTML}{2E7D32}

\newcommand{\R}{\mathbb{R}}
\newcommand{\Rs}{\mathbb{R}^{*}}
\renewcommand{\le}{\leqslant}
\renewcommand{\ge}{\geqslant}
\renewcommand{\approx}{\simeq}
\newcommand{\vect}[1]{\overrightarrow{#1}}
\newcommand{\norme}[1]{\left\lVert #1 \right\rVert}
\newcommand{\intervalle}[2]{\left[#1\,;#2\right]}
\RequirePackage{cancel}
\renewcommand{\CancelColor}{\color{coulExo}}
\newcommand{\fois}[1]{\times{\color{coulExo}#1}}
\newcommand{\barre}[1]{\cancel{#1}}

\tikzset{grille/.style={coulGrille,line width=0.4pt}}
\newcommand{\quadrillage}[4]{\draw[grille,step=1] (#1,#2) grid (#3,#4);}
\tikzset{
  croix/.style={cross out,draw,line width=0.6pt,inner sep=0pt,minimum size=4pt},
  croixdroite/.style={croix,rotate=45,line width=0.8pt},
}
\newcommand{\pt}[3]{\path (#1) node[croix]{} node[#2,font=\small,inner sep=2.5pt]{$#3$};}

\newcommand{\lycee@repere}[4]{%
  \draw[grille,step=1] (#1,#2) grid (#3,#4);
  \draw[-{Stealth[length=2mm]},thick] (#1,0)--({#3+0.4},0) node[below left=-1pt]{$x$};
  \draw[-{Stealth[length=2mm]},thick] (0,#2)--(0,{#4+0.4}) node[below left=-1pt]{$y$};
  \node[below left,font=\scriptsize,inner sep=1.5pt] at (0,0) {$0$};}
\newcommand{\lycee@gradx}[1]{\foreach \k/\l in {#1}{%
  \draw (\k,0.07)--(\k,-0.07) node[below,font=\scriptsize,inner sep=1.5pt]{$\l$};}}
\newcommand{\lycee@grady}[1]{\foreach \k/\l in {#1}{%
  \draw (0.07,\k)--(-0.07,\k) node[left,font=\scriptsize,inner sep=1.5pt]{$\l$};}}
\newcommand{\lycee@intff}[2]{\left[#1\,;#2\right]}
\newcommand{\lycee@intfo}[2]{\left[#1\,;#2\right[}
\newcommand{\lycee@intof}[2]{\left]#1\,;#2\right]}
\newcommand{\lycee@intoo}[2]{\left]#1\,;#2\right[}
\newcommand{\lycee@ens}[1]{\left\{#1\right\}}
\AtBeginDocument{%
  \providecommand{\repere}{\lycee@repere}%
  \providecommand{\gradx}{\lycee@gradx}%
  \providecommand{\grady}{\lycee@grady}%
  \providecommand{\Cf}{\mathcal{C}\sb{\scriptscriptstyle f}}%
  \providecommand{\Cg}{\mathcal{C}\sb{\scriptscriptstyle g}}%
  \providecommand{\intff}{\lycee@intff}%
  \providecommand{\intfo}{\lycee@intfo}%
  \providecommand{\intof}{\lycee@intof}%
  \providecommand{\intoo}{\lycee@intoo}%
  \providecommand{\ens}{\lycee@ens}}

\newlist{questions}{enumerate}{3}
\setlist[questions,1]{label=\arabic*\textdegree),leftmargin=*,itemsep=2pt,topsep=3pt}
\setlist[questions,2]{label=\alph*),leftmargin=*,itemsep=1pt,topsep=2pt}
\setlist[questions,3]{label=\roman*),leftmargin=*,itemsep=1pt,topsep=1pt}
\setlist[itemize]{label=\textbullet,leftmargin=*,itemsep=2pt,topsep=3pt}

\newcommand{\besoinplace}[1]{\par\penalty\@M\begingroup
  \dimen@=#1\relax\dimen@ii=\pagegoal\advance\dimen@ii-\pagetotal
  \ifdim\dimen@>\dimen@ii\ifdim\dimen@ii>\z@\vfil\fi\break\fi\endgroup}

\newcommand{\lycee@pied}{\thepage/\pageref{LastPage}}
\newcommand{\entete}[2]{%
  \pagestyle{fancy}\fancyhf{}%
  \fancyhead[L]{\footnotesize\color{coulPart}\textbf{#1}}%
  \fancyhead[R]{\footnotesize\color{coulPart}\textbf{#2}}%
  \fancyfoot[C]{\footnotesize\lycee@pied}%
  \renewcommand{\headrulewidth}{0.4pt}%
  \renewcommand{\headrule}{\hbox to\headwidth{%
    \color{coulPart!40}\leaders\hrule height \headrulewidth\hfill}}}

\RequirePackage{ccicons}
\newcommand{\lycee@licence}{{\scriptsize\color{gray}%
  \copyright\ Luc Destombe --- \ccbyncsa\ CC BY-NC-SA 4.0}}
\newif\iflycee@licence \lycee@licencetrue
\newcommand{\sanslicence}{\lycee@licencefalse}
\AtBeginDocument{\iflycee@licence\AddToHookNext{shipout/foreground}{%
  \put(0,0){\rlap{\kern\dimexpr1in+\hoffset+\oddsidemargin+\textwidth\relax
    \llap{\raisebox{-\dimexpr1in+\voffset+\topmargin+\headheight+\headsep
      +\textheight+\footskip\relax}{\lycee@licence}}}}}\fi}

\newcommand{\lycee@site}{faire-des-maths.fr}
\newcommand{\lycee@version}{{\scriptsize\color{gray}\lycee@site
  \ifdefined\lyceeversion\ --- Version~\lyceeversion\fi}}
\AtBeginDocument{\iflycee@licence\AddToHookNext{shipout/foreground}{%
  \put(0,0){\rlap{\kern\dimexpr1in+\hoffset+\oddsidemargin\relax
    \raisebox{-\dimexpr1in+\voffset+\topmargin+\headheight+\headsep
      +\textheight+\footskip\relax}{\lycee@version}}}}\fi}

\newcommand{\formulesserrees}{%
  \setlength{\abovedisplayskip}{3pt}\setlength{\belowdisplayskip}{3pt}%
  \setlength{\abovedisplayshortskip}{2pt}\setlength{\belowdisplayshortskip}{3pt}}

\setlength{\parindent}{0pt}

\RequirePackage[implicit=false,hidelinks,pdfencoding=auto,psdextra,bookmarksopen,
  bookmarksopenlevel=1,pdfcreator={},pdfproducer={}]{hyperref}
\RequirePackage{bookmark}
\hypersetup{pdfauthor={Luc Destombe},pdfkeywords={CC BY-NC-SA 4.0}}
\newcounter{lycee@signet}
\newcommand{\signet}[3][]{\stepcounter{lycee@signet}%
  \edef\lycee@dest{\if\relax\detokenize{#1}\relax
    signet.\arabic{lycee@signet}\else#1\fi}%
  \hypertarget{\lycee@dest}{}\bookmark[level=#2,dest=\lycee@dest]{#3}}
\pdfstringdefDisableCommands{%
  \def\R{R}\def\vect#1{#1}\def\Cf{Cf}\def\Cg{Cg}\def\fois#1{×#1}%
  \def\barre#1{#1}\def\intervalle#1#2{[#1;#2]}\def\intff#1#2{[#1;#2]}%
  \def\textsuperscript#1{#1}\def\dfrac#1#2{#1/#2}\def\frac#1#2{#1/#2}%
  \def\vec#1{#1⃗}}%
\newcommand{\res}[1]{#1}
\geometry{top=2cm,bottom=1cm,left=1cm,right=1cm,headheight=14pt,footskip=0.6cm}

\RequirePackage{pgfplots}
\pgfplotsset{compat=1.18}
\RequirePackage{tkz-tab}

\newcounter{partiecours}
\renewcommand{\thepartiecours}{\Roman{partiecours}}
\newcommand{\partiecours}[1]{%
  \refstepcounter{partiecours}%
  \Needspace*{8\baselineskip}%
  \bigskip
  \noindent\begin{tcolorbox}[enhanced,colback=coulPart,colframe=coulPart,
      boxrule=0pt,arc=2pt,left=8pt,right=8pt,top=4pt,bottom=4pt,
      before skip=6pt,after skip=10pt]
    \color{white}\bfseries\large \thepartiecours{} --- #1%
    \signet[partie-\arabic{partiecours}]{1}{\thepartiecours{} --- #1}
  \end{tcolorbox}%
  \addcontentsline{toc}{section}{\thepartiecours{} --- #1}%
}

\newtcolorbox{boitecours}[3][]{%
  enhanced,breakable,
  colback=white,colframe=#2,coltitle=white,
  fonttitle=\bfseries,
  attach boxed title to top left={xshift=8pt,yshift=-\tcboxedtitleheight/2},
  boxed title style={colback=#2,arc=2pt,boxrule=0pt},
  boxrule=0.9pt,arc=2pt,
  left=8pt,right=8pt,top=8pt,bottom=6pt,
  title={#3},#1}

\newcounter{methode}
\newenvironment{definitionbox}[1][Définition]%
  {\begin{boitecours}[phantom={\signet{2}{#1}}]{coulDef}{#1}}{\end{boitecours}}
\newenvironment{proprietebox}[1][Propriété]%
  {\begin{boitecours}[phantom={\signet{2}{#1}}]{coulProp}{#1}}{\end{boitecours}}
\newenvironment{methodebox}[1][Méthode]{\stepcounter{methode}%
  \begin{boitecours}[phantom={\signet[methode-\arabic{methode}]{2}{#1}}]{coulMeth}{#1}}%
  {\end{boitecours}}

\newcommand{\etape}[1]{\par\smallskip\textbf{\color{coulMeth}#1}\par\nobreak\smallskip}
\newcommand{\enonce}[1]{\emph{#1}\par\smallskip}
\newcommand{\reponse}[1]{\par\smallskip
  {\footnotesize\itshape\color{black!55}Réponses : #1}}
\newcommand{\demo}[1]{\textbf{\color{coulMeth}Démonstration.} #1}
\newcommand{\afaire}[1]{%
  \par\smallskip
  \noindent{\small\color{coulExo}$\blacktriangleright$\ \textbf{À faire :}
    fiche d'exercices du chapitre, #1.}\par\smallskip}

\newenvironment{calculs}{\setlength{\jot}{5pt}\@fleqntrue\@mathmargin\retraitcalculs
  \setlength{\abovedisplayskip}{4pt}\setlength{\belowdisplayskip}{4pt}%
  \setlength{\abovedisplayshortskip}{4pt}\setlength{\belowdisplayshortskip}{4pt}%
  \csname alignat*\endcsname{2}}%
  {\csname endalignat*\endcsname}
\newcommand{\rem}[1]{\text{\small\itshape\color{black!60}#1}}
\newcommand{\explique}[2]{\par\smallskip\noindent
  \begin{minipage}[t]{0.57\linewidth}\raggedright #1\end{minipage}\hfill
  \begin{minipage}[t]{0.40\linewidth}\raggedright\leavevmode
    {\small\itshape\color{black!60}#2}\end{minipage}\par}
\newcommand{\cas}[1]{\par\smallskip\noindent\textbf{\color{coulExo}#1}\vadjust{\nobreak}}
\newcommand{\calc}[1]{\par\smallskip\textbf{\color{coulExo}#1}\par\nobreak}

\newtcolorbox{remarquebox}[1][Remarque]{%
  enhanced,breakable,colback=white,colframe=coulRem,
  boxrule=0pt,leftrule=2.5pt,arc=0pt,outer arc=0pt,
  left=8pt,right=6pt,top=5pt,bottom=5pt,
  before upper={\textbf{\color{coulRem}#1}\par\smallskip}}

\newtcolorbox{exemplebox}[1][Exemple]{%
  enhanced,breakable,colback=white,colframe=coulEx,
  boxrule=0pt,leftrule=2.5pt,arc=0pt,outer arc=0pt,
  left=8pt,right=6pt,top=5pt,bottom=5pt,
  before upper={\textbf{\color{coulEx}#1}\par\smallskip}}

\pgfplotsset{
  repere/.style={
    axis lines=middle,
    axis line style={-{Stealth[length=2.4mm]},thick},
    every axis x label/.style={at={(ticklabel* cs:1.0)},anchor=west},
    every axis y label/.style={at={(ticklabel* cs:1.0)},anchor=south},
    label style={font=\footnotesize},
    tick label style={font=\scriptsize},
    grid=both,
    major grid style={grille},
    minor tick num=0,
    tick align=inside,
    clip mode=individual,
  },
}
\tikzset{
  courbe/.style={coulCourbe,line width=1.1pt,smooth},
  courbeB/.style={coulCourbeB,line width=1.1pt,smooth},
  courbeC/.style={coulCourbeC,line width=1.1pt,smooth},
}

\newlist{etapes}{enumerate}{1}
\setlist[etapes]{label=\textbf{\arabic*.},leftmargin=*,itemsep=1pt,topsep=2pt}

\newlist{expressions}{enumerate}{1}
\setlist[expressions]{label={\Alph*\ $=$},leftmargin=*,itemsep=3pt,topsep=3pt}

\newcounter{exo}
\renewcommand{\theexo}{\ifnum\value{exo}<10 0\fi\arabic{exo}}
\newtcolorbox{exercice}[1][]{%
  enhanced,breakable,
  colback=white,colframe=coulExo,
  boxrule=0.9pt,arc=2pt,
  left=8pt,right=8pt,top=8pt,bottom=6pt,
  coltitle=white,fonttitle=\bfseries,
  attach boxed title to top left={xshift=8pt,yshift=-\tcboxedtitleheight/2},
  boxed title style={colback=coulExo,arc=2pt,boxrule=0pt},
  phantom={\signet{2}{Exercice \arabic{exo}}},
  title={Exercice \theexo\ifx\relax#1\relax\else{} --- #1\fi}}
\newenvironment{exo}[1][\relax]{\refstepcounter{exo}\begin{exercice}[#1]}{\end{exercice}}

  \newenvironment{correction}{\par}{\par}
  \newcommand{\autableau}{}
  \newcommand{\etapeexemples}{\etape{Exemples corrigés}}
  \newcommand{\etapetableau}[1]{\etape{#1}}
  \newenvironment{exempletableau}[1][Exemple]%
    {\begin{exemplebox}[#1]}{\end{exemplebox}}

\RequirePackage{comment}
\excludecomment{correctionavj}

\newcommand{\coursEntete}{}
\newcommand{\coursNiveau}{}
\newcommand{\coursTitre}{}
\newcommand{\coursSurTitre}{}
\newcommand{\coursSousTitre}{}

\newcommand{\enteteCours}[1]{\renewcommand{\coursEntete}{#1}}
\newcommand{\chapitrecours}[2]{\enteteCours{Chapitre #1 --- #2}}
\newcommand{\niveaucours}[1]{\renewcommand{\coursNiveau}{#1}}
\newcommand{\titrecours}[3][]{%
  \renewcommand{\coursTitre}{#2}%
  \renewcommand{\coursSurTitre}{#1}%
  \renewcommand{\coursSousTitre}{#3}}

\pagestyle{fancy}
\fancyhf{}
\fancyhead[L]{\footnotesize\color{coulPart}\textbf{\coursEntete}}
\fancyhead[R]{\footnotesize\color{coulPart}\textbf{\coursNiveau}}
\fancyfoot[C]{\footnotesize\thepage}
\renewcommand{\headrulewidth}{0.4pt}
\renewcommand{\headrule}{\hbox to\headwidth{\color{coulPart!40}\leaders\hrule height \headrulewidth\hfill}}

\setlength{\parskip}{2pt}

\AtBeginDocument{%
  \begin{center}
    \begin{tcolorbox}[enhanced,width=0.72\linewidth,colback=white,colframe=coulPart,
        boxrule=1.6pt,arc=3pt,halign=center,top=10pt,bottom=10pt]
      \ifx\coursSurTitre\empty
        {\Huge\bfseries\color{coulPart} \coursTitre}\\[4pt]
      \else
        {\Huge\bfseries\color{coulPart} \coursTitre}\\[3pt]
        {\Large\bfseries\color{coulPart} \coursSurTitre}\\[5pt]
      \fi
      {\normalsize \coursSousTitre}%
      
    \end{tcolorbox}
  \end{center}%
}

\usetikzlibrary{decorations.text}
\definecolor{lycee@metal}{HTML}{1F4E79}
\definecolor{lycee@manche}{HTML}{8A5A2B}
\definecolor{lycee@etincelle}{HTML}{F28C28}
\newcommand{\lycee@pluriel}[2]{#1 #2\ifnum#1>1 s\fi}
\newcommand{\lycee@sceau}{%
  \edef\lycee@bas{\lycee@pluriel{\lyceerelectures}{relecture} et
    \lycee@pluriel{\lyceerodages}{rodage}}%
  \begin{tikzpicture}[line join=round]
    \fill[coulPart!10,draw=coulPart,line width=.8pt] (0:2.2)
      \foreach \i in {1,...,71}{ -- (\i*5:{isodd(\i)?2.1:2.2}) } -- cycle;
    \draw[coulPart,line width=1.2pt] (0,0) circle (1.93);
    \draw[coulPart,line width=.4pt] (0,0) circle (1.85);
    \draw[coulPart,line width=.4pt] (0,0) circle (1.2);
    \path[decorate,decoration={text along path,text align=center,
        text={|\fontsize{9.5}{11}\selectfont\bfseries\color{coulPart}|TEST{\'E} AVEC MES {\'E}L{\`E}VES}}]
      (180:1.45) arc (180:0:1.45);
    \path[decorate,decoration={text along path,text align=center,
        text/.expanded={|\noexpand\fontsize{8.5}{10}\noexpand\selectfont\noexpand\color{coulPart}|\lycee@bas}}]
      (195:1.68) arc (195:345:1.68);
    \begin{scope}[x=.25mm,y=-.25mm]
      \begin{scope}[shift={(-32,-40)}]%
        \fill[lycee@metal] (23,31) rectangle (41,49);
        \draw[lycee@etincelle,line width=.5pt,line cap=round]
          (20,30)--(15,27) (19,34)--(14,35) (44,30)--(49,27) (45,34)--(50,35);
        \fill[lycee@manche,rounded corners=.3mm] (38,24.6) rectangle (60,27.6);
        \fill[lycee@metal,rounded corners=.4mm] (25,20) rectangle (39,31);
        \fill[lycee@etincelle,rounded corners=.2mm] (25,20) rectangle (28,31);
      \end{scope}
    \end{scope}
    \ifdefined\lyceesceauplusieurs\def\lycee@annee{Ann{\'e}es~:}\else
      \def\lycee@annee{Ann{\'e}e}\fi
    \path[decorate,decoration={text along path,text align=center,
        text/.expanded={|\noexpand\fontsize{8}{10}\noexpand\selectfont\noexpand\color{coulPart}|%
          \unexpanded\expandafter{\lycee@annee} \unexpanded\expandafter{\lyceesceauannees}}}]
      (95:1.08) arc (95:445:1.08);
  \end{tikzpicture}}
\AtEndDocument{\ifdefined\lyceesceauannees
  \par\besoinplace{5cm}\vspace{1.5em}%
  \begin{center}\lycee@sceau\end{center}\fi}
\makeatother

\chapitrecours{9}{Équations de droites}
\niveaucours{Seconde}
\titrecours{ÉQUATIONS DE DROITES}{Cours --- Seconde}

\begin{document}

\ShorthandsOff
\tikzset{
  vecteur/.style={-{Stealth[length=2.6mm]},line width=1.1pt,coulCourbe},
  vecteurB/.style={-{Stealth[length=2.6mm]},line width=1.1pt,coulCourbeB},
  droite/.style={coulCourbe,line width=1.1pt},
  droiteB/.style={coulCourbeB,line width=1.1pt},
  droiteC/.style={coulCourbeC,line width=1.1pt},
  aide/.style={dashed,black!45,line width=0.6pt},
  fig/.style={line width=0.8pt,black!75},
}
\ShorthandsOn
\newcommand{\coord}[2]{\begin{pmatrix}#1\\#2\end{pmatrix}}

\begin{remarquebox}[Comment utiliser ce cours]
Chaque partie commence par ce qu'il faut \textbf{savoir} (définitions, propriétés), puis
vient ce qu'il faut \textbf{savoir faire} : une méthode, avec des
\textbf{exemples corrigés} faits en classe, puis des exercices
\og\textbf{À vous de jouer}\fg{} à chercher seul. Le repère est orthonormé dans tout le
chapitre.

\end{remarquebox}

\partiecours{Vecteur directeur d'une droite}

\begin{definitionbox}[Définition 1 --- Vecteur directeur]
Un \textbf{vecteur directeur} d'une droite $(d)$ est un vecteur non nul qui a la même
direction que $(d)$. Si $A$ et $B$ sont deux points distincts de $(d)$, le vecteur
$\vect{AB}$ est un vecteur directeur de $(d)$.

\smallskip
Une droite a une infinité de vecteurs directeurs, tous colinéaires : si $\vec{u}$ en est
un, $2\vec{u}$, $-\vec{u}$\ldots{} en sont d'autres.
\end{definitionbox}

\begin{methodebox}[Méthode 1 --- Lire un vecteur directeur, tracer une droite]
Pour \textbf{lire}, on cherche deux endroits où la droite passe exactement par un
\textbf{nœud} du quadrillage (le croisement de deux lignes) ; on les marque d'une croix et
on compte le déplacement de l'un à l'autre. Pour \textbf{tracer} la droite
qui passe par un point $A$ avec un vecteur directeur $\vec{u}$, on fait le déplacement de
$\vec{u}$ depuis $A$ : on obtient un second point.
\etapeexemples
\begin{minipage}[t]{0.52\linewidth}\raggedright
\begin{questions}[label=\alph*)]
  \item Lire un vecteur directeur de $(d_{1})$, puis un autre.
  \item Lire un vecteur directeur de $(d_{2})$.
  \item Tracer la droite $(d_{3})$ qui passe par $A(-1\,;-2)$ et de vecteur directeur
        $\vec{u}\coord{1}{3}$.
\end{questions}
\end{minipage}\hfill
\begin{minipage}[t]{0.45\linewidth}
\vspace{0pt}\centering
\begin{tikzpicture}[x=0.5cm,y=0.5cm]
  \repere{-4}{-3}{5}{4}
  \gradx{-3,-2,-1,1,2,3,4} \grady{-2,-1,1,2,3}
  \draw[droite] (-4,-1.5) -- (5,3);
  \draw[droiteB] (-3.5,4) -- (5,-1.6667);
  \node[coulCourbe,font=\footnotesize,above left] at (4.6,2.8) {$(d_{1})$};
  \node[coulCourbeB,font=\footnotesize,right] at (-3.3,4) {$(d_{2})$};
\end{tikzpicture}
\end{minipage}
\begin{correction}
\cas{a)}\par\nopagebreak
\begin{minipage}[t]{0.58\linewidth}
\explique{$(d_{1})$ passe exactement par les nœuds $(-3\,;-1)$ et $(-1\,;0)$, marqués d'une
croix : de l'un à l'autre, $2$ carreaux vers la droite, $1$ vers le haut.}{ailleurs, elle
coupe les carreaux entre deux nœuds}
\explique{$\res{\vec{u}_{1}\coord{2}{1}}$ est un vecteur directeur de $(d_{1})$ ;
$\res{2\vec{u}_{1}\coord{4}{2}}$ en est un autre.}{tout vecteur colinéaire convient}
\cas{b)}\par\nopagebreak
\explique{$(d_{2})$ passe par les nœuds $(-2\,;3)$ et $(1\,;1)$ :
$\res{\vec{u}_{2}\coord{3}{-2}}$ est un vecteur directeur.}{$3$ vers la droite, $2$ vers le
bas}
\end{minipage}\hfill
\begin{minipage}[t]{0.4\linewidth}
\vspace{0pt}\centering
\begin{tikzpicture}[x=0.45cm,y=0.45cm]
  \repere{-4}{-3}{5}{4}
  \draw[droite] (-4,-1.5) -- (5,3);
  \draw[droiteB] (-3.5,4) -- (5,-1.6667);
  \draw[vecteur,black] (-3,-1) -- node[midway,below right,font=\footnotesize]{$\vec{u}_{1}$} (-1,0);
  \draw[vecteur,black] (-2,3) -- node[midway,above right,font=\footnotesize]{$\vec{u}_{2}$} (1,1);
  \path (-3,-1) node[croix]{} (-1,0) node[croix]{} (-2,3) node[croix]{} (1,1) node[croix]{};
\end{tikzpicture}
\end{minipage}
\cas{c)}\par\nopagebreak
\begin{minipage}[t]{0.58\linewidth}
\explique{Depuis $A(-1\,;-2)$, $1$ carreau vers la droite et $3$ vers le haut : on arrive
en $(0\,;1)$.}{déplacement de $\vec{u}$}
\explique{On trace la droite qui passe par $A$ et par ce point.}{deux points suffisent ;
on peut continuer jusqu'à $(1\,;4)$ pour vérifier}
\end{minipage}\hfill
\begin{minipage}[t]{0.4\linewidth}
\vspace{0pt}\centering
\begin{tikzpicture}[x=0.4cm,y=0.4cm]
  \repere{-3}{-3}{3}{4}
  \draw[droiteC] (-1.3333,-3) -- (1,4);
  \draw[vecteur,black] (-1,-2) -- node[midway,left,font=\footnotesize]{$\vec{u}$} (0,1);
  \pt{-1,-2}{left}{A}
  \node[coulCourbeC,font=\footnotesize,right] at (1,3.6) {$(d_{3})$};
\end{tikzpicture}
\end{minipage}
\end{correction}

\etape{À vous de jouer}
\begin{minipage}[t]{0.52\linewidth}\raggedright
\begin{questions}[label=\alph*)]
  \item Lire un vecteur directeur de $(d_{4})$, puis un autre.
  \item Lire un vecteur directeur de $(d_{5})$.
  \item Tracer la droite $(d_{6})$ qui passe par $B(-2\,;1)$ et de vecteur directeur
        $\vec{v}\coord{3}{1}$.
\end{questions}
\end{minipage}\hfill
\begin{minipage}[t]{0.45\linewidth}
\vspace{0pt}\centering
\begin{tikzpicture}[x=0.5cm,y=0.5cm]
  \repere{-4}{-3}{5}{4}
  \gradx{-3,-2,-1,1,2,3,4} \grady{-2,-1,1,2,3}
  \draw[droite] (-4,1.75) -- (5,4);
  \draw[droiteB] (-1.75,4) -- (3.5,-3);
  \node[coulCourbe,font=\footnotesize,below] at (-3.2,1.7) {$(d_{4})$};
  \node[coulCourbeB,font=\footnotesize,right] at (3.3,-2.6) {$(d_{5})$};
\end{tikzpicture}
\end{minipage}
\reponse{a) $\coord{4}{1}$, puis par exemple $\coord{8}{2}$ ; b)~$\coord{3}{-4}$ ;
c)~la droite passe par $B$, $(1\,;2)$ et $(4\,;3)$.}
\end{methodebox}

\afaire{exercices 1 à 4}

\partiecours{Équation cartésienne d'une droite}

\begin{proprietebox}[Théorème 1 --- Équation cartésienne]
\begin{itemize}
  \item Toute droite $(d)$ a une équation de la forme $ax+by+c=0$, où $a$ et $b$ ne
        sont pas nuls en même temps : un point $M(x\,;y)$ est sur $(d)$ quand ses coordonnées
        vérifient l'égalité. C'est une \textbf{équation cartésienne} de $(d)$.
  \item Le vecteur $\vec{u}\coord{-b}{a}$ est un vecteur directeur de $(d)$.
  \item Réciproquement (admis), l'ensemble des points $M(x\,;y)$ tels que $ax+by+c=0$ ($a$
        et $b$ n'étant pas nuls en même temps) est une droite.
\end{itemize}
\end{proprietebox}

\textbf{Exemple, puis démonstration}\autableau
\begin{correction}
\explique{Exemple : la droite qui passe par $A(1\,;2)$, de vecteur directeur
$\vec{u}\coord{3}{1}$. Un point $M(x\,;y)$ est sur la droite quand $\vect{AM}\coord{x-1}{y-2}$
et $\vec{u}$ sont colinéaires.}{chapitre 8 : colinéarité par le déterminant}
\begin{calculs}
\det(\vect{AM}\,;\vec{u}) &= (x-1)\times 1-(y-2)\times 3 &\qquad& \rem{$xy'-yx'$}\\
\det(\vect{AM}\,;\vec{u}) &= x-1-3y+6 && \rem{on développe}\\
\det(\vect{AM}\,;\vec{u}) &= x-3y+5 && \rem{on réduit}
\end{calculs}
\explique{Une équation de la droite est $x-3y+5=0$ ; on retrouve $\vec{u}\coord{-b}{a}$ avec
$a=1$ et $b=-3$.}{déterminant nul}
\explique{Démonstration : $(d)$ passe par $A(x_{A}\,;y_{A})$ ; on nomme $-b$ et $a$ les
coordonnées d'un vecteur directeur $\vec{u}$. $M(x\,;y)$ est sur $(d)$ quand
$\vect{AM}\coord{x-x_{A}}{y-y_{A}}$ et $\vec{u}\coord{-b}{a}$ sont colinéaires.}{même
démarche que l'exemple}
\begin{calculs}
\det(\vect{AM}\,;\vec{u}) &= (x-x_{A})\times a-(y-y_{A})\times(-b) &\qquad& \rem{$xy'-yx'$}\\
\det(\vect{AM}\,;\vec{u}) &= ax-ax_{A}+by-by_{A} && \rem{on développe}\\
\det(\vect{AM}\,;\vec{u}) &= ax+by+c && \rem{en posant $c=-ax_{A}-by_{A}$}
\end{calculs}
\explique{Conclusion : $M$ est sur $(d)$ quand $ax+by+c=0$.}{déterminant nul}
\end{correction}

\begin{methodebox}[Méthode 2 --- Exploiter une équation cartésienne]
On lit $a$, $b$ et $c$ (avec leurs signes). Un point est sur la droite quand ses
coordonnées vérifient l'équation. Pour tracer, il suffit d'un point et d'un vecteur
directeur, ou de deux points.
\etapeexemples
La droite $(d)$ a pour équation $2x+3y-6=0$.
\begin{questions}[label=\alph*)]
  \item Donner un vecteur directeur de $(d)$.
  \item Les points $A(3\,;0)$ et $C(1\,;2)$ sont-ils sur $(d)$ ?
  \item Calculer l'ordonnée du point de $(d)$ d'abscisse $1$.
  \item Tracer la droite $(d)$.
\end{questions}
\begin{correction}
\cas{a)}\par\nopagebreak
\explique{$a=2$ et $b=3$, donc $\res{\vec{u}\coord{-3}{2}}$ est un vecteur
directeur.}{$\vec{u}\coord{-b}{a}$ : on change le signe de $b$, qu'on écrit en haut}
\cas{b)}
\begin{calculs}
2x_{A}+3y_{A}-6 &= 2\times 3+3\times 0-6 &\qquad& \rem{on remplace $x$ par $3$, $y$ par $0$}\\
2x_{A}+3y_{A}-6 &= 0 &&
\end{calculs}
\explique{Les coordonnées de $A$ vérifient l'équation : $A$ \res{est sur $(d)$}.}{égalité
vraie}
\begin{calculs}
2x_{C}+3y_{C}-6 &= 2\times 1+3\times 2-6 &\qquad& \rem{on remplace $x$ par $1$, $y$ par $2$}\\
2x_{C}+3y_{C}-6 &= 2 &&
\end{calculs}
\explique{$2\neq 0$ : $C$ \res{n'est pas sur $(d)$}.}{égalité fausse}
\cas{c)}
\begin{calculs}
2\times 1+3y-6 &= 0 &\qquad& \rem{on remplace $x$ par $1$}\\
3y-4 &= 0 && \rem{on réduit}\\
3y &= 4 && \rem{on ajoute $4$}\\
y &= \res{\frac{4}{3}} && \rem{on divise par $3$}
\end{calculs}
\cas{d)}\par\nopagebreak
\begin{minipage}[t]{0.58\linewidth}
\explique{$A(3\,;0)$ est sur $(d)$ (question b).}{un point déjà connu}
\explique{Depuis $A$, le déplacement de $\vec{u}\coord{-3}{2}$ mène en $(0\,;2)$. On trace
la droite par ces deux points.}{vecteur directeur de la question a}
\end{minipage}\hfill
\begin{minipage}[t]{0.4\linewidth}
\vspace{0pt}\centering
\begin{tikzpicture}[x=0.4cm,y=0.4cm]
  \repere{-4}{-2}{5}{5}
  \gradx{-3,3} \grady{2,4}
  \draw[droite] (-4,4.6667) -- (5,-1.3333);
  \draw[vecteur,black] (3,0) -- node[midway,below left,font=\footnotesize]{$\vec{u}$} (0,2);
  \pt{3,0}{above right}{A}
  \node[coulCourbe,font=\footnotesize,above right] at (-3,4) {$(d)$};
\end{tikzpicture}
\end{minipage}
\end{correction}

\etape{À vous de jouer}
La droite $(d')$ a pour équation $x+3y-4=0$.
\begin{questions}[label=\alph*)]
  \item Donner un vecteur directeur de $(d')$.
  \item Les points $E(1\,;1)$ et $F(-2\,;1)$ sont-ils sur $(d')$ ?
  \item Calculer l'ordonnée du point de $(d')$ d'abscisse $2$.
  \item Tracer la droite $(d')$.
\end{questions}
\reponse{a) $\coord{-3}{1}$ ; b)~$E$ est sur $(d')$, $F$ non ; c)~$\frac{2}{3}$ ;
d)~par $E(1\,;1)$ et $(-2\,;2)$.}
\end{methodebox}

\afaire{exercices 5 à 8}

\partiecours{Trouver une équation cartésienne}

\begin{remarquebox}[Lien avec le théorème 1]
Un vecteur directeur $\vec{u}\coord{-b}{a}$ donne $a$ (en bas) et $b$ (en haut, avec le
signe changé). Il reste $c$ : on le trouve en écrivant qu'un point connu de la droite
vérifie l'équation. Une droite a plusieurs équations cartésiennes : en multipliant
$ax+by+c=0$ par un même nombre non nul, on en obtient une autre.
\end{remarquebox}

\begin{methodebox}[Méthode 3 --- Trouver une équation cartésienne]
\etapeexemples
\begin{questions}[label=\alph*)]
  \item Trouver une équation cartésienne de la droite $(d)$ qui passe par $A(2\,;-1)$ et
        de vecteur directeur $\vec{u}\coord{4}{3}$.
  \item Trouver une équation cartésienne de la droite $(BC)$, avec $B(-1\,;2)$ et
        $C(3\,;4)$.
\end{questions}
\begin{correction}
\cas{a)}
\explique{$\vec{u}\coord{-b}{a}=\coord{4}{3}$, donc $a=3$ et $b=-4$ : $(d)$ a une équation
$3x-4y+c=0$.}{$-b=4$, donc $b=-4$}
\begin{calculs}
3\times 2-4\times(-1)+c &= 0 &\qquad& \rem{$A$ est sur $(d)$ : on remplace $x$ et $y$}\\
10+c &= 0 && \rem{$6+4=10$}\\
c &= -10 &&
\end{calculs}
\explique{Une équation cartésienne de $(d)$ est $\res{3x-4y-10=0}$.}{contrôle avec $A$ :
$6+4-10=0$}
\cas{b)}
\explique{$\vect{BC}\coord{4}{2}$ est un vecteur directeur : $a=2$ et $b=-4$. $(BC)$ a une
équation $2x-4y+c=0$.}{$\coord{3-(-1)}{4-2}$}
\begin{calculs}
2\times(-1)-4\times 2+c &= 0 &\qquad& \rem{$B$ est sur la droite}\\
-10+c &= 0 && \rem{$-2-8=-10$}\\
c &= 10 &&
\end{calculs}
\explique{Une équation de $(BC)$ est $\res{2x-4y+10=0}$, ou $x-2y+5=0$ en divisant par
$2$.}{contrôle avec $C$ : $3-8+5=0$}
\end{correction}

\etape{À vous de jouer}
\begin{questions}[label=\alph*)]
  \item Trouver une équation cartésienne de la droite $(d)$ qui passe par $A(-3\,;1)$ et
        de vecteur directeur $\vec{u}\coord{2}{5}$.
  \item Trouver une équation cartésienne de la droite $(EF)$, avec $E(1\,;-2)$ et
        $F(4\,;4)$.
\end{questions}
\reponse{a) $5x-2y+17=0$ ; b)~$6x-3y-12=0$, ou $2x-y-4=0$.}
\end{methodebox}

\afaire{exercices 9 à 12}

\partiecours{Équation réduite ; passer de l'une à l'autre}

\begin{proprietebox}[Propriété 1 --- Équation réduite]
Si $b\neq 0$, on isole $y$ dans $ax+by+c=0$ : la droite a une équation unique de la forme
\[
  y=mx+p,
\]
son \textbf{équation réduite}. C'est la représentation graphique de la fonction affine
${f(x)=mx+p}$, notée ${f(x)=ax+b}$ au chapitre 5 : on change de lettres, car $a$ et $b$
servent déjà dans l'équation cartésienne. Ainsi :
\begin{itemize}
  \item $m$ est la \textbf{pente}, ou \textbf{coefficient directeur} : quand $x$ augmente
        de $1$, $y$ augmente de $m$ ; le vecteur $\vec{u}\coord{1}{m}$ est un vecteur
        directeur ;
  \item $p$ est l'\textbf{ordonnée à l'origine} : la droite passe par le point $(0\,;p)$ ;
  \item si $A$ et $B$ sont deux points de la droite :
        $m=\dfrac{y_{B}-y_{A}}{x_{B}-x_{A}}$ (propriété 3 du chapitre 5).
\end{itemize}
Si $b=0$, la droite est parallèle à l'axe des ordonnées (partie VI).
\end{proprietebox}

\begin{methodebox}[Méthode 4 --- Passer d'une équation à l'autre]
Pour la réduite, on isole $y$ comme dans une équation. Pour revenir à une équation
cartésienne, on fait tout passer du même côté (et on multiplie pour chasser un
dénominateur).
\etapeexemples
\begin{questions}[label=\alph*)]
  \item Trouver l'équation réduite de la droite d'équation $4x+2y-6=0$.
  \item Trouver l'équation réduite de la droite d'équation $2x-3y+6=0$, puis sa pente et
        son ordonnée à l'origine.
  \item Trouver une équation cartésienne de la droite d'équation $y=\frac{1}{3}x-2$.
\end{questions}
\begin{correction}
\cas{a)}
\begin{calculs}
4x+2y-6 &= 0 &\qquad& \rem{on recopie}\\
2y &= -4x+6 && \rem{on enlève $4x$, on ajoute $6$}\\
y &= \res{-2x+3} && \rem{on divise tout par $2$}
\end{calculs}
\cas{b)}
\begin{calculs}
2x-3y+6 &= 0 &\qquad& \rem{on recopie}\\
-3y &= -2x-6 && \rem{on enlève $2x$ et $6$}\\
y &= \res{\frac{2}{3}x+2} && \rem{on divise tout par $-3$ : les signes changent}
\end{calculs}
\explique{Pente \res{$\frac{2}{3}$}, ordonnée à l'origine \res{$2$}.}{$m=\frac23$ et $p=2$}
\cas{c)}
\begin{calculs}
y &= \frac{1}{3}x-2 &\qquad& \rem{on recopie}\\
3y &= x-6 && \rem{on multiplie tout par $3$}\\
0 &= x-3y-6 && \rem{on enlève $3y$ des deux côtés}
\end{calculs}
\explique{Une équation cartésienne est $\res{x-3y-6=0}$.}{contrôle : $(0\,;-2)$ la vérifie}
\end{correction}

\etape{À vous de jouer}
\begin{questions}[label=\alph*)]
  \item Trouver l'équation réduite de la droite d'équation $6x+2y-8=0$.
  \item Trouver l'équation réduite de la droite d'équation $3x-5y+10=0$, puis sa pente et
        son ordonnée à l'origine.
  \item Trouver une équation cartésienne de la droite d'équation $y=\frac{2}{3}x+1$.
\end{questions}
\reponse{a) $y=-3x+4$ ; b)~$y=\frac{3}{5}x+2$, pente $\frac35$, ordonnée à l'origine $2$ ;
c)~$2x-3y+3=0$.}
\end{methodebox}

\afaire{exercices 13 à 17}

\partiecours{Droites parallèles}

\begin{proprietebox}[Propriété 2 --- Droites parallèles, droites sécantes]
\begin{itemize}
  \item Deux droites sont \textbf{parallèles} quand leurs vecteurs directeurs sont
        colinéaires (déterminant nul, chapitre 8).
  \item Pour $(d)$ : $ax+by+c=0$ et $(d')$ : $a'x+b'y+c'=0$, on compare
        $\vec{u}\coord{-b}{a}$ et $\vec{u}'\coord{-b'}{a'}$.
  \item Pour $(d)$ : $y=mx+p$ et $(d')$ : $y=m'x+p'$, les droites sont parallèles quand
        elles ont la \textbf{même pente} : $m=m'$. En effet,
        $\det\left(\coord{1}{m}\,;\coord{1}{m'}\right)=m'-m$, nul quand $m=m'$.
  \item Deux droites qui ne sont pas parallèles sont \textbf{sécantes} : elles se coupent
        en un point.
\end{itemize}
\end{proprietebox}

\begin{methodebox}[Méthode 5 --- Reconnaître des droites parallèles]
On lit un vecteur directeur de chaque droite, puis on calcule leur déterminant : nul, les
droites sont parallèles ; non nul, elles sont sécantes.
\etapeexemples
Pour chaque couple de droites, dire si les droites sont parallèles ou sécantes.
\begin{questions}[label=\alph*)]
  \item $(d_{1})$ : $2x-3y+1=0$ et $(d_{2})$ : $-4x+6y+5=0$.
  \item $(d_{3})$ : $x+2y-3=0$ et $(d_{4})$ : $2x-y+1=0$.
  \item $(d_{5})$ : $x-3y+1=0$ et $(d_{6})$ : $\frac{1}{3}x-y+2=0$.
\end{questions}
\begin{correction}
\cas{a)}
\explique{$\vec{u}_{1}\coord{3}{2}$ dirige $(d_{1})$ et $\vec{u}_{2}\coord{-6}{-4}$ dirige
$(d_{2})$.}{$\vec{u}\coord{-b}{a}$}
\begin{calculs}
\det(\vec{u}_{1}\,;\vec{u}_{2}) &= 3\times(-4)-2\times(-6) &\qquad& \rem{$xy'-yx'$}\\
\det(\vec{u}_{1}\,;\vec{u}_{2}) &= -12+12 &&\\
\det(\vec{u}_{1}\,;\vec{u}_{2}) &= 0 &&
\end{calculs}
\explique{Conclusion : $(d_{1})$ et $(d_{2})$ sont \res{parallèles}.}{vecteurs directeurs
colinéaires}
\cas{b)}
\explique{$\vec{u}_{3}\coord{-2}{1}$ et $\vec{u}_{4}\coord{1}{2}$.}{$\vec{u}\coord{-b}{a}$ :
$-b=-(-1)=1$ pour $(d_{4})$}
\begin{calculs}
\det(\vec{u}_{3}\,;\vec{u}_{4}) &= -2\times 2-1\times 1 &\qquad& \rem{$xy'-yx'$}\\
\det(\vec{u}_{3}\,;\vec{u}_{4}) &= -5 &&
\end{calculs}
\explique{Conclusion : $(d_{3})$ et $(d_{4})$ sont \res{sécantes}.}{$-5\neq 0$}
\cas{c)}
\explique{$\vec{u}_{5}\coord{3}{1}$ et $\vec{u}_{6}\coord{1}{\frac13}$.}{$a=\frac13$ et $b=-1$
pour $(d_{6})$}
\begin{calculs}
\det(\vec{u}_{5}\,;\vec{u}_{6}) &= 3\times\frac{1}{3}-1\times 1 &\qquad& \rem{$xy'-yx'$}\\
\det(\vec{u}_{5}\,;\vec{u}_{6}) &= 0 && \rem{$1-1$}
\end{calculs}
\explique{Conclusion : $(d_{5})$ et $(d_{6})$ sont \res{parallèles}.}{vecteurs directeurs
colinéaires}
\end{correction}

\etape{À vous de jouer}
Pour chaque couple de droites, dire si les droites sont parallèles ou sécantes.
\begin{questions}[label=\alph*)]
  \item $(d_{1})$ : $3x-y+2=0$ et $(d_{2})$ : $-6x+2y-1=0$.
  \item $(d_{3})$ : $2x+y-4=0$ et $(d_{4})$ : $x-3y+1=0$.
  \item $(d_{5})$ : $x+3y-2=0$ et $(d_{6})$ : $\frac{1}{3}x+y+5=0$.
\end{questions}
\reponse{a) parallèles ; b)~sécantes ; c)~parallèles.}
\end{methodebox}

\afaire{exercices 18 et 19}

\begin{methodebox}[Méthode 6 --- Trouver une équation de la parallèle qui passe par un point]
Une parallèle a le même vecteur directeur : elle garde les mêmes $a$ et $b$, seul $c$
change. On le trouve avec le point.
\etapeexemples
\begin{questions}[label=\alph*)]
  \item Trouver une équation cartésienne de la parallèle à $(d)$ : $x-2y+3=0$ qui passe
        par $A(1\,;5)$.
  \item Trouver une équation cartésienne de la parallèle à $(d')$ : $2x+5y-1=0$ qui passe
        par $B(-1\,;3)$.
\end{questions}
\begin{correction}
\cas{a)}
\explique{La parallèle a une équation $x-2y+c=0$.}{mêmes $a$ et $b$ : même vecteur
directeur}
\begin{calculs}
1-2\times 5+c &= 0 &\qquad& \rem{$A(1\,;5)$ est sur la parallèle}\\
-9+c &= 0 && \rem{$1-10=-9$}\\
c &= 9 &&
\end{calculs}
\explique{Une équation de la parallèle est $\res{x-2y+9=0}$.}{contrôle : $1-10+9=0$}
\cas{b)}
\explique{La parallèle a une équation $2x+5y+c=0$.}{mêmes $a$ et $b$}
\begin{calculs}
2\times(-1)+5\times 3+c &= 0 &\qquad& \rem{$B$ est sur la parallèle}\\
13+c &= 0 && \rem{$-2+15=13$}\\
c &= -13 &&
\end{calculs}
\explique{Une équation de la parallèle est $\res{2x+5y-13=0}$.}{contrôle avec $B$}
\end{correction}

\etape{À vous de jouer}
\begin{questions}[label=\alph*)]
  \item Trouver une équation cartésienne de la parallèle à $(d)$ : $3x+y-7=0$ qui passe
        par $A(3\,;-1)$.
  \item Trouver une équation cartésienne de la parallèle à $(d')$ : $3x-4y+2=0$ qui passe
        par $B(2\,;-1)$.
\end{questions}
\reponse{a) $3x+y-8=0$ ; b)~$3x-4y-10=0$.}
\end{methodebox}

\afaire{exercices 20 et 21}

\partiecours{Droites horizontales et verticales}

\begin{proprietebox}[Propriété 3 --- Droites parallèles aux axes]
\begin{minipage}[t]{0.6\linewidth}\raggedright
Dans une équation cartésienne $ax+by+c=0$ :
\begin{itemize}
  \item si $a=0$, l'équation $by+c=0$ se ramène à $y=k$ : tous les points ont la même
        ordonnée, la droite est \textbf{horizontale} (parallèle à l'axe des
        abscisses) ; $\vec{u}\coord{-b}{0}$ la dirige ;
  \item si $b=0$, l'équation $ax+c=0$ se ramène à $x=k$ : tous les points ont la même
        abscisse, la droite est \textbf{verticale} (parallèle à l'axe des ordonnées) ;
        $\vec{u}\coord{0}{a}$ la dirige.
\end{itemize}
\end{minipage}\hfill
\begin{minipage}[t]{0.36\linewidth}
\vspace{0pt}\centering
\begin{tikzpicture}[x=0.45cm,y=0.45cm]
  \repere{-3}{-2}{4}{4}
  \gradx{-2,2} \grady{-1,3}
  \draw[droite] (-3,3) -- (4,3);
  \draw[droiteB] (2,-2) -- (2,4);
  \node[coulCourbe,font=\footnotesize,above] at (-1.8,3) {$y=3$};
  \node[coulCourbeB,font=\footnotesize,right] at (2,-1.4) {$x=2$};
\end{tikzpicture}
\end{minipage}
\end{proprietebox}

\begin{methodebox}[Méthode 7 --- Reconnaître et tracer une droite horizontale ou verticale]
\etapeexemples
\begin{questions}[label=\alph*)]
  \item Tracer les droites $(d_{1})$ : $2y+2=0$ et $(d_{2})$ : $x+2=0$.
  \item Trouver une équation cartésienne de la droite $(AB)$, avec $A(4\,;-1)$ et
        $B(4\,;3)$.
  \item Trouver une équation cartésienne de la droite $(CD)$, avec $C(-2\,;5)$ et
        $D(3\,;5)$.
  \item Quelle est la nature de la droite d'équation $2x+6=0$ ?
\end{questions}
\begin{correction}
\cas{a)}\par\nopagebreak
\begin{minipage}[t]{0.58\linewidth}
\explique{$(d_{1})$ : $2y+2=0$ donne $y=-1$, une droite horizontale.}{$a=0$ ; tous les
points d'ordonnée $-1$}
\explique{$(d_{2})$ : $x+2=0$ donne $x=-2$, une droite verticale.}{$b=0$ ; tous les points
d'abscisse $-2$}
\end{minipage}\hfill
\begin{minipage}[t]{0.4\linewidth}
\vspace{0pt}\centering
\begin{tikzpicture}[x=0.4cm,y=0.4cm]
  \repere{-4}{-3}{3}{2}
  \gradx{-2} \grady{-1}
  \draw[droite] (-4,-1) -- (3,-1);
  \draw[droiteB] (-2,-3) -- (-2,2);
  \node[coulCourbe,font=\footnotesize,below] at (2,-1) {$(d_{1})$};
  \node[coulCourbeB,font=\footnotesize,left] at (-2,1.5) {$(d_{2})$};
\end{tikzpicture}
\end{minipage}
\cas{b)}
\explique{$\vect{AB}\coord{0}{4}$ : $a=4$ et $b=0$, d'où $4x+c=0$.}{méthode 3 ;
$-b=0$}
\begin{calculs}
4\times 4+c &= 0 &\qquad& \rem{$A$ est sur la droite}\\
c &= -16 &&
\end{calculs}
\explique{$\res{4x-16=0}$, soit $x=4$ : $(AB)$ est verticale.}{$A$ et $B$ ont la même
abscisse}
\cas{c)}
\explique{$\vect{CD}\coord{5}{0}$ : $a=0$ et $b=-5$, d'où $-5y+c=0$.}{méthode 3 ;
$-b=5$}
\begin{calculs}
-5\times 5+c &= 0 &\qquad& \rem{$C$ est sur la droite}\\
c &= 25 &&
\end{calculs}
\explique{$\res{-5y+25=0}$, soit $y=5$ : $(CD)$ est horizontale.}{$C$ et $D$ ont la même
ordonnée}
\cas{d)}
\begin{calculs}
2x+6 &= 0 &\qquad& \rem{pas de $y$ : $b=0$}\\
x &= -3 && \rem{on isole $x$}
\end{calculs}
\explique{C'est \res{la droite verticale d'équation $x=-3$}.}{tous les points
d'abscisse $-3$}
\end{correction}

\etape{À vous de jouer}
\begin{questions}[label=\alph*)]
  \item Tracer les droites $(d_{3})$ : $3y-6=0$ et $(d_{4})$ : $x-1=0$.
  \item Trouver une équation cartésienne de la droite $(EF)$, avec $E(-3\,;2)$ et
        $F(-3\,;-4)$.
  \item Trouver une équation cartésienne de la droite $(GH)$, avec $G(1\,;-1)$ et
        $H(5\,;-1)$.
  \item Quelle est la nature de la droite d'équation $4x-8=0$ ?
\end{questions}
\reponse{b) $-6x-18=0$, soit $x=-3$ ; c)~$-4y-4=0$, soit $y=-1$ ; d)~la droite verticale
d'équation $x=2$.}
\end{methodebox}

\afaire{exercices 22 à 24}

\partiecours{Point d'intersection de deux droites sécantes}

\begin{proprietebox}[Propriété 4 --- Point d'intersection]
Deux droites sécantes d'équations réduites $y=mx+p$ et $y=m'x+p'$ se coupent en un point
$M(x\,;y)$ dont les coordonnées vérifient les deux équations :
\begin{itemize}
  \item son abscisse $x$ est la solution de l'équation $mx+p=m'x+p'$ ;
  \item son ordonnée s'obtient en remplaçant $x$ dans l'une des deux équations.
\end{itemize}
\end{proprietebox}

\begin{methodebox}[Méthode 8 --- Calculer les coordonnées du point d'intersection]
On égale les deux expressions de $y$ et on résout ; on calcule $y$ avec l'une des
équations, et on contrôle avec l'autre.
\etapeexemples
\begin{questions}[label=\alph*)]
  \item Calculer les coordonnées du point d'intersection de $(d_{1})$ : $y=2x-3$ et
        $(d_{2})$ : $y=-x+3$.
  \item Même question avec $(d_{3})$ : $y=2x+1$ et $(d_{4})$ : $y=-x+3$.
\end{questions}
\begin{correction}
\cas{a)}\par\nopagebreak
\begin{minipage}[t]{0.58\linewidth}\raggedright
\begin{calculs}
2x-3 &= -x+3 &\qquad& \rem{même $y$}\\
3x &= 6 && \rem{on ajoute $x$ et $3$}\\
x &= 2 && \rem{on divise par $3$}
\end{calculs}
\begin{calculs}
y &= 2\times 2-3 &\qquad& \rem{dans $(d_{1})$}\\
y &= 1 &&
\end{calculs}
\explique{Le point d'intersection est $\res{(2\,;1)}$.}{contrôle dans $(d_{2})$ :
$-2+3=1$}
\end{minipage}\hfill
\begin{minipage}[t]{0.38\linewidth}
\vspace{0pt}\centering
\begin{tikzpicture}[x=0.45cm,y=0.45cm]
  \repere{-1}{-3}{4}{4}
  \gradx{1,2,3} \grady{-2,-1,1,2,3}
  \draw[droite] (0,-3) -- (3.5,4);
  \draw[droiteB] (-1,4) -- (4,-1);
  \draw[aide] (2,0) -- (2,1) -- (0,1);
  \path (2,1) node[croix]{};
  \node[coulCourbe,font=\footnotesize,left] at (3.3,3.6) {$(d_{1})$};
  \node[coulCourbeB,font=\footnotesize,below left] at (4,-1) {$(d_{2})$};
\end{tikzpicture}
\end{minipage}
\cas{b)}
\begin{calculs}
2x+1 &= -x+3 &\qquad& \rem{même $y$}\\
3x &= 2 && \rem{on ajoute $x$, on enlève $1$}\\
x &= \frac{2}{3} && \rem{on divise par $3$}
\end{calculs}
\begin{calculs}
y &= 2\times\frac{2}{3}+1 &\qquad& \rem{dans $(d_{3})$}\\
y &= \frac{4}{3}+\frac{1\fois{3}}{1\fois{3}} && \rem{même dénominateur}\\
y &= \frac{7}{3} &&
\end{calculs}
\explique{Le point d'intersection est
$\res{\left(\frac{2}{3}\,;\frac{7}{3}\right)}$.}{contrôle dans $(d_{4})$ :
$-\frac23+\frac93=\frac73$}
\end{correction}

\etape{À vous de jouer}
\begin{questions}[label=\alph*)]
  \item Calculer les coordonnées du point d'intersection de $(d_{1})$ : $y=3x-4$ et
        $(d_{2})$ : $y=-x+4$.
  \item Même question avec $(d_{3})$ : $y=-2x+1$ et $(d_{4})$ : $y=x+3$.
\end{questions}
\reponse{a) $(2\,;2)$ ; b)~$\left(-\frac{2}{3}\,;\frac{7}{3}\right)$.}
\end{methodebox}

\afaire{exercices 25 à 28}

\partiecours{Problèmes de géométrie repérée}

\begin{remarquebox}[Pourquoi un repère ?]
Une figure sans coordonnées (un carré, un rectangle\ldots) peut se traiter par le calcul :
on choisit un repère orthonormé adapté (un sommet à l'origine, deux côtés sur les axes),
on écrit les coordonnées des points, puis chaque question devient un calcul du chapitre 8
ou de ce chapitre : parallélisme et alignement par le déterminant, milieu par les
moyennes, longueurs par la norme, intersection par les équations réduites.
\end{remarquebox}

\begin{methodebox}[Méthode 9 --- Résoudre un problème de géométrie avec un repère]
\etapeexemples
\begin{minipage}[t]{0.62\linewidth}\raggedright
$ABCD$ est un carré de côté $6$~cm ; $I$, $J$ et $K$ sont les milieux de $[AB]$, $[BC]$
et $[CD]$. On se place dans le repère orthonormé d'origine $A$, d'axes $(AB)$ et $(AD)$,
d'unité $1$~cm.
\begin{questions}[label=\alph*)]
  \item Donner les coordonnées des points $A$, $B$, $C$, $D$, $I$, $J$ et $K$.
  \item Démontrer que les droites $(DI)$ et $(BK)$ sont parallèles.
  \item Démontrer que le triangle $IJK$ est rectangle isocèle.
  \item Calculer les coordonnées du point $L$, intersection de $(AJ)$ et $(DI)$.
\end{questions}
\end{minipage}\hfill
\begin{minipage}[t]{0.35\linewidth}
\vspace{0pt}\centering
\begin{tikzpicture}[x=0.45cm,y=0.45cm]
  \draw[fig] (0,0) -- (6,0) -- (6,6) -- (0,6) -- cycle;
  \draw[aide] (0,6) -- (3,0) (0,0) -- (6,3) (6,0) -- (3,6);
  \pt{0,0}{below left}{A} \pt{6,0}{below right}{B} \pt{6,6}{above right}{C}
  \pt{0,6}{above left}{D} \pt{3,0}{below}{I} \pt{6,3}{right}{J} \pt{3,6}{above}{K}
\end{tikzpicture}
\end{minipage}
\begin{correction}
\cas{a)}\par\nopagebreak
\explique{$\res{A(0\,;0)}$, $\res{B(6\,;0)}$, $\res{C(6\,;6)}$ et
$\res{D(0\,;6)}$.}{côté $6$ ; repère orthonormé : angle droit en $A$, même unité sur les
deux axes}
\explique{$\res{I(3\,;0)}$, $\res{J(6\,;3)}$ et $\res{K(3\,;6)}$.}{moyennes des coordonnées
des extrémités}
\cas{b)}
\explique{$\vect{DI}\coord{3}{-6}$ et $\vect{BK}\coord{-3}{6}$.}{extrémité moins origine}
\begin{calculs}
\det(\vect{DI}\,;\vect{BK}) &= 3\times 6-(-6)\times(-3) &\qquad& \rem{$xy'-yx'$}\\
\det(\vect{DI}\,;\vect{BK}) &= 0 && \rem{$18-18$}
\end{calculs}
\explique{Conclusion : $(DI)$ et $(BK)$ sont \res{parallèles}.}{vecteurs colinéaires}
\cas{c)}
\explique{$\vect{IJ}\coord{3}{3}$, $\vect{JK}\coord{-3}{3}$ et
$\vect{IK}\coord{0}{6}$.}{extrémité moins origine}
\begin{calculs}
IJ^{2} &= 3^{2}+3^{2} &\qquad& \rem{carrés des longueurs}\\
IJ^{2} &= 18 &&
\end{calculs}
\begin{calculs}
JK^{2} &= (-3)^{2}+3^{2} &\qquad& \rem{même calcul}\\
JK^{2} &= 18 &&
\end{calculs}
\begin{calculs}
IK^{2} &= 0^{2}+6^{2} &\qquad& \rem{le plus long côté}\\
IK^{2} &= 36 &&
\end{calculs}
\explique{$IJ=JK$ et $IJ^{2}+JK^{2}=36=IK^{2}$.}{isocèle ; réciproque de Pythagore}
\explique{Conclusion : $IJK$ est \res{rectangle isocèle en $J$}.}{on rassemble}
\cas{d)}
\explique{$(AJ)$ passe par l'origine et par $J(6\,;3)$ : pente $\frac{3}{6}=\frac12$,
équation $y=\frac{1}{2}x$.}{$p=0$ : la droite passe par $A(0\,;0)$}
\explique{$(DI)$ : pente $\frac{0-6}{3-0}=-2$ et ordonnée à l'origine $6$ :
$y=-2x+6$.}{$D(0\,;6)$ est sur l'axe des ordonnées}
\begin{calculs}
\frac{1}{2}x &= -2x+6 &\qquad& \rem{même $y$}\\
\frac{5}{2}x &= 6 && \rem{on ajoute $2x$ : $\frac12+2=\frac52$}\\
x &= \frac{12}{5} && \rem{on multiplie par $\frac25$}
\end{calculs}
\begin{calculs}
y &= \frac{1}{2}\times\frac{12}{5} &\qquad& \rem{dans $(AJ)$}\\
y &= \frac{6}{5} &&
\end{calculs}
\explique{Donc $\res{L\left(\frac{12}{5}\,;\frac{6}{5}\right)}$.}{contrôle dans $(DI)$ :
$-\frac{24}{5}+\frac{30}{5}=\frac65$}
\end{correction}

\etape{À vous de jouer}
\begin{minipage}[t]{0.62\linewidth}\raggedright
$ABCD$ est un rectangle avec $AB=8$~cm et $AD=4$~cm ; $I$, $J$ et $K$ sont les milieux de
$[AB]$, $[BC]$ et $[CD]$. On se place dans le repère orthonormé d'origine $A$, d'axes
$(AB)$ et $(AD)$, d'unité $1$~cm.
\begin{questions}[label=\alph*)]
  \item Donner les coordonnées des points $A$, $B$, $C$, $D$, $I$, $J$ et $K$.
  \item Démontrer que les droites $(DI)$ et $(BK)$ sont parallèles.
  \item Démontrer que le triangle $DIC$ est rectangle isocèle.
  \item Calculer les coordonnées du point $L$, intersection de $(AJ)$ et $(DI)$.
\end{questions}
\end{minipage}\hfill
\begin{minipage}[t]{0.35\linewidth}
\vspace{0pt}\centering
\begin{tikzpicture}[x=0.35cm,y=0.35cm]
  \draw[fig] (0,0) -- (8,0) -- (8,4) -- (0,4) -- cycle;
  \draw[aide] (0,4) -- (4,0) (0,0) -- (8,2) (8,0) -- (4,4);
  \pt{0,0}{below left}{A} \pt{8,0}{below right}{B} \pt{8,4}{above right}{C}
  \pt{0,4}{above left}{D} \pt{4,0}{below}{I} \pt{8,2}{right}{J} \pt{4,4}{above}{K}
\end{tikzpicture}
\end{minipage}
\reponse{a) $A(0\,;0)$, $B(8\,;0)$, $C(8\,;4)$, $D(0\,;4)$, $I(4\,;0)$, $J(8\,;2)$,
$K(4\,;4)$ ; b)~$\vect{DI}\coord{4}{-4}$ et $\vect{BK}\coord{-4}{4}$, déterminant $0$ ;
c)~$DI^{2}=IC^{2}=32$ et $DC^{2}=64$ ; d)~$L\left(\frac{16}{5}\,;\frac{4}{5}\right)$.}
\end{methodebox}

\afaire{exercices 29 à 31}

\end{document}
